
\documentclass[10pt, fleqn]{article}

% Page Layout
\usepackage[a4paper, margin=0.5in]{geometry}
\usepackage{parskip}
\usepackage{multicol}
\usepackage{titlesec}
\usepackage{enumitem}
\usepackage{float}
\usepackage[font=small, skip=4pt]{caption}

% Mathematics
\usepackage{amsmath, amssymb, amsfonts, mathtools, bm, mathrsfs}

% Figures and Tables
\usepackage{graphicx}
\usepackage{tabularx, booktabs}

% Additional Formatting
\usepackage{siunitx}
\usepackage{soul}
\usepackage[hidelinks]{hyperref}

% US Customary Units
\DeclareSIUnit{\inch}{in}
\DeclareSIUnit{\foot}{ft}
\DeclareSIUnit{\yard}{yd}
\DeclareSIUnit{\mile}{mi}

\DeclareSIUnit{\slug}{slug}
\DeclareSIUnit{\lbm}{lbm}
\DeclareSIUnit{\lbf}{lbf}

\DeclareSIUnit{\psi}{psi}
\DeclareSIUnit{\psf}{psf}

\DeclareSIUnit{\gallon}{gal}
\DeclareSIUnit{\quart}{qt}
\DeclareSIUnit{\fahrenheit}{F}
\DeclareSIUnit{\rankine}{R}

\sisetup{
    per-mode = fraction
}

% Equation Spacing
\setlength{\abovedisplayskip}{6pt}
\setlength{\belowdisplayskip}{6pt}
\setlength{\abovedisplayshortskip}{4pt}
\setlength{\belowdisplayshortskip}{4pt}
\setlength{\mathindent}{0pt}

% Figure and Table Spacing
\setlength{\textfloatsep}{8pt plus 2pt minus 2pt}
\setlength{\intextsep}{6pt plus 2pt minus 2pt}
\setlength{\floatsep}{8pt plus 2pt minus 2pt}

% Section Formatting
\titleformat{\section}{\large\bfseries}{}{0em}{}
\titleformat{\subsection}{\normalsize\bfseries}{}{0em}{}

\begin{document}

\begin{center}
    \Large \textbf{Feedback Controls Formula Sheet} \\
    \normalsize Exam 1
\end{center}
    \section*{System Modeling and Analysis}
    \vspace{-1em}
\begin{multicols*}{2}\noindent
\textbf{State Space Representation}\\
\[\begin{aligned}
\dot{\underline{x}} &= [A]{\underline{x}} + [B]\underline{u} \\
\underline{y} &= [C]{\underline{x}} + [D]\underline{u}
\end{aligned}\]
\underline{Mass-Spring-Damper System:}
\[m\ddot{x}+ b\dot{x} + kx = f(t)\]
Define the state variables to "kill the dots":
\[x = x,\quad v=\dot{x}\]
Rewrite the equation of motion as a system of first order differential equations: $m\dot{v}+ bv + kx = f(t)$
\begin{eqnarray*}
\dot{x} &=& v \\
\dot{v} &=& -\frac{k}{m}x - \frac{b}{m}v + \frac{1}{m}f(t)
\end{eqnarray*}
In state space form,
\[
\underbrace{
\begin{bmatrix}
\dot{x}\\
\dot{v}
\end{bmatrix}
}_{\dot{\underline{x}}}
=
\underbrace{
\begin{bmatrix}
0 & 1\\
-\frac{k}{m} & -\frac{b}{m}
\end{bmatrix}
}_{[A]}
\underbrace{
\begin{bmatrix}
x\\
v
\end{bmatrix}
}_{\underline{x}}
+
\underbrace{
\begin{bmatrix}
0\\
\frac{1}{m}
\end{bmatrix}
}_{[B]}
\underbrace{f(t)}_{\underline{u}}
\]
\[
\underbrace{y}_{\underline{y}}
=
\underbrace{
\begin{bmatrix}
1 & 0
\end{bmatrix}
}_{[C]}
\underbrace{
\begin{bmatrix}
x\\
v
\end{bmatrix}
}_{\underline{x}}
+
\underbrace{0}_{[D]}
\underbrace{f(t)}_{\underline{u}}
\]
\textbf{Transfer Function Matrix from State Space}
\[\frac{\underline{Y}}{\underline{U}} = [C][sI - A]^{-1}[B] + [D]\]
\underline{Recall:} We can invert a 2x2 matrix using the formula:
\[
A =
\begin{pmatrix}
a & b \\
c & d
\end{pmatrix}\]
\[A^{-1} = \frac{1}{ad - bc}
\begin{pmatrix}
d & -b \\
-c & a
\end{pmatrix},
\quad \text{if } ad - bc \neq 0.\]
\textbf{Linearization (Jacobian Method)}\\
Nonlinear pendulum: torque input $T = k_\tau \sin(u_1)$ and force input $F = k_F \sin(u_2)$, where $u_1$ and $u_2$ are independent inputs:
\begin{center}
    \includegraphics[width=0.3\textwidth]{C:/Users/zsong/Desktop/LaTex/Fall 2026 Notes/Feedback Controls/pendulum+.png}
\end{center}
The equation of motion is given by:
\[\sum \tau = I \ddot{\theta} \implies m\ell^2 \ddot{\theta} = -mg\ell\sin\theta + k_\tau \sin(u_1) + \ell k_F \sin(u_2)\]
Define $\omega = \dot{\theta}$ and the state vector:
\[\underline{x} = \begin{pmatrix} \theta \\ \omega \end{pmatrix},\quad \underline{u} = \begin{pmatrix} u_1 \\ u_2 \end{pmatrix}\]
Nonlinear state-space representation:
\[\dot{\underline{x}} = \begin{pmatrix} \omega \\ -\frac{g}{\ell}\sin\theta + \frac{k_\tau}{m\ell^2}\sin(u_1) + \frac{k_F}{m\ell}\sin(u_2) \end{pmatrix} \iff \underline{\dot{x}} = \underline{f}(\underline{x}, \underline{u})\]
\[y=\begin{pmatrix}
    1 & 0
\end{pmatrix} \underline{x} \quad \iff \quad y = \underline{h}(\underline{x}, \underline{u})\quad \text{(output $y = \theta$)}\]
Linearization procedure:
\begin{itemize}
    \item[I.] Define operating point $(\underline{\bar{x}}, \underline{\bar{u}})$ \& incremental variables:
    \[\underline{x} = \underline{\bar{x}} + \underline{\hat{x}},\quad \underline{u} = \underline{\bar{u}} + \underline{\hat{u}},\quad \underline{y} = \underline{\bar{y}} + \underline{\hat{y}}\]
    Linearize the nonlinear state equation:
    \[\underline{f}(\underline{x},\underline{u}) \approx \underline{f}(\underline{\bar{x}}, \underline{\bar{u}}) + \frac{\partial \underline{f}}{\partial \underline{x}}\bigg|_{(\underline{\bar{x}}, \underline{\bar{u}})} \underline{\hat{x}} + \frac{\partial \underline{f}}{\partial \underline{u}}\bigg|_{(\underline{\bar{x}}, \underline{\bar{u}})} \underline{\hat{u}}\]
    \begin{itemize}
        \item State Jacobian $[A]$:
        \[\renewcommand{\arraystretch}{1.8}
        [A] = \left.\begin{pmatrix}
\frac{\partial f_1}{\partial \theta} & \frac{\partial f_1}{\partial \omega} \\
\frac{\partial f_2}{\partial \theta} & \frac{\partial f_2}{\partial \omega}
\end{pmatrix}\right|_{(\underline{\bar{x}}, \underline{\bar{u}})}=
\begin{pmatrix}
0 & 1 \\
-\frac{g}{\ell}\cos\bar{\theta} & 0
\end{pmatrix}\]
        \item Input Jacobian $[B]$:
\[
\renewcommand{\arraystretch}{1.8}
[B] =
\left.
\begin{pmatrix}
\dfrac{\partial f_1}{\partial u_1} & \dfrac{\partial f_1}{\partial u_2} \\
\dfrac{\partial f_2}{\partial u_1} & \dfrac{\partial f_2}{\partial u_2}
\end{pmatrix}
\right|_{(\underline{\bar{x}}, \underline{\bar{u}})} =
\begin{pmatrix}
0 & 0 \\
\dfrac{k_{\tau}}{m\ell^2}\cos\bar{u}_1 &
\dfrac{k_F}{m\ell}\cos\bar{u}_2
\end{pmatrix}\]
    \end{itemize}
    \item [II.] Find operating point: $\dot{\underline{\bar{x}}} = \underline{f}(\underline{{\bar{x}}},\underline{\bar{u}}) = 0$
    \[\dot{\underline{x}}\big|_{(\underline{\bar{x}}, \underline{\bar{u}})} = \begin{pmatrix} \bar{\omega} \\ -\frac{g}{\ell}\sin\bar{\theta} + \frac{k_\tau}{m\ell^2}\sin(\bar{u_1}) + \frac{k_F}{m\ell}\sin(\bar{u_2}) \end{pmatrix} = \begin{pmatrix} 0 \\ 0 \end{pmatrix}\]
    Solving these conditions define $(\underline{\bar{x}}, \underline{\bar{u}})$. Leave it symbolic for now.
    \item [III.] Substitute linearized terms into state equation and apply equilibrium condition:
    \[\big(\underline{\dot{\bar{x}}}+\underline{\dot{\hat{x}}}\big)\approx \underline{f}(\underline{{\bar{x}}},\underline{\bar{u}})+ A \underline{\hat{x}}\ + B \underline{\hat{u}}\]
    Equilibrium condition $\underline{f}(\underline{{\bar{x}}},\underline{\bar{u}}) = 0$ gives:
    \[\underline{\dot{\hat{x}}} = A \underline{\hat{x}} + B \underline{\hat{u}}\]
\[
\resizebox{0.95\columnwidth}{!}{$\displaystyle
\renewcommand{\arraystretch}{2}
\boxed{
\underbrace{
\begin{pmatrix}
\dot{\hat{\theta}} \\
\dot{\hat{\omega}}
\end{pmatrix}
}_{\dot{\hat{x}}}
=
\underbrace{
\begin{pmatrix}
0 & 1 \\
-\dfrac{g}{\ell}\cos\bar{\theta} & 0
\end{pmatrix}
}_{A}
\underbrace{
\begin{pmatrix}
\hat{\theta} \\
\hat{\omega}
\end{pmatrix}
}_{\hat{x}}
+
\underbrace{
\begin{pmatrix}
0 & 0 \\
\dfrac{k_{\tau}}{m\ell^2}\cos\bar{u}_1 &
\dfrac{k_F}{m\ell}\cos\bar{u}_2
\end{pmatrix}
}_{B}
\underbrace{
\begin{pmatrix}
\hat{u}_1 \\
\hat{u}_2
\end{pmatrix}
}_{\hat{u}}
}
$}
\]
\end{itemize}
Output can be linearized in a similar manner:
\[
C = \left.\frac{\partial h}{\partial x}\right|_{(\underline{\bar{x}}, \underline{\bar{u}})}
=\left.\begin{pmatrix} \frac{\partial h}{\partial \theta} & \frac{\partial h}{\partial \omega} \end{pmatrix}\right|_{(\underline{\bar{x}}, \underline{\bar{u}})}
= \begin{pmatrix} 1 & 0 \end{pmatrix}
\]
\[
D = \left.\frac{\partial h}{\partial u}\right|_{(\underline{\bar{x}}, \underline{\bar{u}})}
=\left.\begin{pmatrix} \frac{\partial h}{\partial u_1} & \frac{\partial h}{\partial u_2} \end{pmatrix}\right|_{(\underline{\bar{x}}, \underline{\bar{u}})}
= \begin{pmatrix} 0 & 0 \end{pmatrix}
\]
\[
\implies \hat{y} = C\hat{x} + D\hat{u}
= \begin{pmatrix} 1 & 0 \end{pmatrix}\hat{x}
\]
\underline{Note:} For a scalar nonlinear system, linearization is simply a first-order Taylor expansion about the operating point:
\[
f(x)\approx f(\bar{x})+\left.\frac{df}{dx}\right|_{x=\bar{x}}(x-\bar{x})
\]
\newpage

\textbf{Laplace Transforms}

\begin{center}
\renewcommand{\arraystretch}{1.15}
\begin{tabular}{@{}ll@{}}
\midrule
$\delta(t)$ & $1$ \\
$1$ & $1/s$ \\
$t$ & $1/s^2$ \\
$e^{at}$ & $1/(s-a)$ \\
$\dot{f}(t)$ & $sF(s)-f(0^-)$ \\
$\ddot{f}(t)$ & $s^2F(s)-sf(0^-)-\dot{f}(0^-)$ \\
\bottomrule
\end{tabular}
\end{center}
\textbf{Initial and Final Value Theorems:}
\[
\boxed{f(0^+)=\lim_{s\to\infty}sF(s)}
\qquad
\boxed{f(\infty)=\lim_{s\to 0}sF(s)}
\]
\textbf{Transfer Functions and Block Diagrams}\\
\[G(s) = \frac{Y(s)}{U(s)} = \frac{\text{"Forward Gain"}}{1+\text{"Loop Gain"}}\]
Series Blocks: $G(s) = G_1(s)G_2(s)$ \\
Parallel Blocks: $G(s) = G_1(s)+G_2(s)$ 
\section*{System Behavior and Root Locus}
    \vspace{-1em}
\textbf{DC Gain and Steady State Error}
\[K_{DC} = \lim_{s \to 0} sY(s) = \lim_{s \to 0} G(s) \quad \text{(For unit step input)}\]
\[e_{ss} = \lim_{s \to 0} sE(s),\quad e(t) = r(t) - y(t)\]
\textbf{First Order System}
\[\tau \dot{y} + y = K u\leftrightarrow G(s) = \frac{K}{\tau s + 1}\]
\[y(t) = K(1 - e^{-t/\tau}) \quad \text{(For unit step input)}\]
\textbf{Second Order System}
\[\ddot{y} + 2\zeta\omega_n\dot{y} + \omega_n^2 y = \omega_n^2 u\leftrightarrow G(s) = \frac{\omega_n^2}{s^2 + 2\zeta\omega_n s + \omega_n^2}\]
\[\sigma = \zeta\omega_n \quad \text{(exponential damping rate)}\]
\[\omega_d = \omega_n\sqrt{1 - \zeta^2} \quad \text{(damped frequency)}\]
Poles:
$s_{1,2} = -\zeta\omega_n \pm \omega_n\sqrt{\zeta^2 - 1} = -\sigma \pm \hat{\imath}\omega_d $
\[\phi = \cos^{-1}\left(\frac{\sigma}{|s_{1,2}|}\right) = \cos^{-1}\left(\frac{\zeta\omega_n}{\omega_n}\right) = \cos^{-1}(\zeta)\]
\begin{center}
    \includegraphics[width=0.25\textwidth]{C:/Users/zsong/Desktop/LaTex/Fall 2026 Notes/Feedback Controls/2nd order poles.png}
\end{center}
For unit step input, the time response is given by:
 \[\implies y(t) = 1 - \frac{1}{\sqrt{1-\zeta^2}} e^{-\zeta\omega_n t} \sin(\omega_d t + \phi)\]
 \newcolumn
\begin{center}
\renewcommand{\arraystretch}{1.6}
\begin{tabular}{@{}ll@{}}
\midrule
Undamped ($\zeta=0$)
& $A\cos(\omega_n t)+B\sin(\omega_n t)$ \\

Underdamped ($0<\zeta<1$)
& $e^{-\zeta\omega_n t}[A\cos(\omega_d t)+B\sin(\omega_d t)]$ \\

Critically Damped ($\zeta=1$)
& $(C_1+C_2t)e^{-\omega_n t}$ \\

Overdamped ($\zeta>1$)
& $C_1e^{s_1t}+C_2e^{s_2t}$ \\
\bottomrule
\end{tabular}
\end{center}
\textbf{Time-Domain Specs}\\
For second order underdamped system given unit step input:
\vspace{-2em}
\begin{center}
    \includegraphics[width=0.8\linewidth]{C:/Users/zsong/Desktop/LaTex/Fall 2026 Notes/Feedback Controls/Time specs.png}
\end{center}
\begin{center}
\renewcommand{\arraystretch}{1.7}
\begin{tabular}{@{}ll@{}}
\midrule
Rise Time (0--100\%)
& $\displaystyle t_r=\frac{\pi-\phi}{\omega_d}
=\frac{\pi-\cos^{-1}(\zeta)}{\omega_n\sqrt{1-\zeta^2}}$ \\

Rise Time (10--90\%)
& $\displaystyle t_r\approx\frac{1.8}{\omega_n}$ \\

Peak Time
& $\displaystyle t_p=\frac{\pi}{\omega_d}
=\frac{\pi}{\omega_n\sqrt{1-\zeta^2}}$ \\

Max. Overshoot
& $\displaystyle M_p=e^{-\frac{\pi\zeta}{\sqrt{1-\zeta^2}}}\times100\%$ \\


& $\implies \displaystyle \zeta=\frac{-\ln(M_p)}
{\sqrt{\pi^2+[\ln(M_p)]^2}}$ \\

Settling Time (2\%)
& $\displaystyle t_s\approx\frac{4}{\zeta\omega_n}$ \\

Settling Time (5\%)
& $\displaystyle t_s\approx\frac{3.6}{\zeta\omega_n}$ \\
\bottomrule
\end{tabular}
\end{center}
\textbf{Evans Root Locus Rules}\\
\[1 + KL_0(s) = 0,\quad K \geq 0\quad \text{(Evans Form)}\]
\[\boxed{|K L_0(s)| = 1\implies K = \frac{1}{|L_0(s)|} \quad \text{(magnitude condition)}}\]
\[ \boxed{\angle{L_0(s)} = (2\ell+1)\pi,\quad \ell \in \mathbb{Z}\quad \text{(phase condition)}}\]
Rule 0: The root locus is symmetric about the real axis. \\\\
Rule 1: Branches of the root locus start at the poles of $L_0(s)$ and end at the zeros of $L_0(s)$ or infinity (Asymptotes). \\\\
Rule 2: A point $s$ on the real axis is part of the root locus if
$k_0 > 0 : \quad N_R \text{ is odd},\quad 
k_0 < 0 : \quad N_R \text{ is even}$
where $N_R$ is the number of poles and zeros to the right of $s$. \\\\
Rule 3: All asymptotes intersect at the centroid $\alpha$ with direction $\theta_q$, let relative degree $r = $  number of poles $-$ number of zeros:
{\scriptsize
\[
\boxed{\alpha=\frac{\sum p_i-\sum z_j}{r}}
\quad
\boxed{\theta_q=
\begin{cases}
(2q+1)\pi/r,&k_0>0\\
2q\pi/r,&k_0<0
\end{cases}
\quad q=0,\ldots,r-1}
\]
}










\end{multicols*}









\end{document}
